% ============================================================
% Chapter 8: Simple Linear Regression
% ============================================================
\chapter{Simple Linear Regression}
\label{ch:simple-regression}

\begin{center}
\textit{``Regression is the most useful and most misused tool in statistics.''}
\end{center}

\section{Motivating Example}
Data on living area (m$^2$) and sale price (kEUR) for 15 apartments. Goal: model $\text{Price}=f(\text{Area})$.

\section{The Model}

\begin{definition}[Simple Linear Regression Model]
\index{linear regression!simple}
$Y_i = \beta_0 + \beta_1 x_i + \varepsilon_i$, $i=1,\ldots,n$, with $\E[\varepsilon_i]=0$, $\Var(\varepsilon_i)=\sigma^2$ (homoscedasticity), $\Cov(\varepsilon_i,\varepsilon_j)=0$, and optionally $\varepsilon_i\sim\mathcal{N}(0,\sigma^2)$.
\end{definition}

\section{OLS Estimation}

\begin{theorem}[OLS Estimators]
\index{OLS}
\[
\hat{\beta}_1 = \frac{S_{xy}}{S_{xx}} = \frac{\sum(x_i-\bar{x})(Y_i-\bar{Y})}{\sum(x_i-\bar{x})^2}, \qquad \hat{\beta}_0 = \bar{Y}-\hat{\beta}_1\bar{x}.
\]
\end{theorem}

\begin{proof}
Minimize $Q=\sum(Y_i-\beta_0-\beta_1 x_i)^2$. Setting $\partial Q/\partial\beta_0=0$ gives $\beta_0=\bar{Y}-\beta_1\bar{x}$. Substituting and solving $\partial Q/\partial\beta_1=0$ yields the result.
\end{proof}

\begin{theorem}[Gauss--Markov]
\index{Gauss--Markov}
Under assumptions (1)--(3), OLS estimators are BLUE (Best Linear Unbiased Estimators).
\end{theorem}

\section{Variance Decomposition}

\begin{theorem}
$\underbrace{\sum(Y_i-\bar{Y})^2}_{\text{SST}} = \underbrace{\sum(\hat{Y}_i-\bar{Y})^2}_{\text{SSR}} + \underbrace{\sum(Y_i-\hat{Y}_i)^2}_{\text{SSE}}$.
\end{theorem}

\begin{definition}[$R^2$]
\index{R-squared@$R^2$}
$R^2 = \text{SSR}/\text{SST} = 1-\text{SSE}/\text{SST} = r_{xy}^2$ in simple regression.
\end{definition}

\section{Inference under Normality}

$T_1 = \hat{\beta}_1/\mathrm{SE}(\hat{\beta}_1) \sim t(n-2)$. Test $H_0:\beta_1=0$.

$F = \text{SSR}/1 \;/\; \text{SSE}/(n-2) \sim F(1,n-2)$. In simple regression, $F=T_1^2$.

\subsection{Prediction}

CI for $\E[Y|x_0]$: $\hat{Y}_0\pm t_{\alpha/2,n-2}\hat{\sigma}\sqrt{1/n+(x_0-\bar{x})^2/S_{xx}}$.

PI for $Y_{\text{new}}$: add 1 under the square root.

\section{Residual Diagnostics}
\index{residuals}
Check: (1) residuals vs fitted (linearity, homoscedasticity), (2) QQ-plot (normality), (3) residuals vs order (autocorrelation).

\section{Python Implementation}

\begin{python}
\begin{minted}{python}
import numpy as np
from scipy import stats
import statsmodels.api as sm
import matplotlib.pyplot as plt

area = np.array([25,30,35,40,45,50,55,60,65,70,75,80,90,100,120])
price = np.array([95,110,125,140,148,160,175,185,195,210,220,240,260,290,350])

slope, intercept, r, p, se = stats.linregress(area, price)
print(f"beta_1={slope:.4f} (p={p:.6f}), R^2={r**2:.4f}")

X = sm.add_constant(area)
model = sm.OLS(price, X).fit()
print(model.summary())

fig, axes = plt.subplots(2, 2, figsize=(12, 10))
axes[0,0].scatter(area, price); axes[0,0].plot(area, model.fittedvalues, 'r-')
axes[0,0].set_title('Regression line')
axes[0,1].scatter(model.fittedvalues, model.resid)
axes[0,1].axhline(0, color='r', ls='--'); axes[0,1].set_title('Residuals vs Fitted')
stats.probplot(model.resid, plot=axes[1,0]); axes[1,0].set_title('QQ-plot')
axes[1,1].hist(model.resid, bins=6, edgecolor='black')
axes[1,1].set_title('Residual distribution')
plt.tight_layout(); plt.savefig("ch08_regression.pdf"); plt.show()
\end{minted}
\end{python}

\section{Exercises}

\begin{exercise}[$\star$]
Compute $\hat{\beta}_0$, $\hat{\beta}_1$, $R^2$ by hand for $x=(1,2,3,4,5)$, $y=(2.1,3.8,6.2,7.9,10.3)$.
\end{exercise}

\begin{exercise}[$\star\star$]
Prove Gauss--Markov for simple regression: among $\hat{\beta}_1=\sum c_i Y_i$ linear unbiased, OLS has minimum variance.
\end{exercise}

\begin{exercise}[$\star\star\star$ -- Project]
Collect real data, fit a simple regression, test $\beta_1=0$, analyze residuals, construct prediction intervals.
\end{exercise}

\begin{keyformulas}
\begin{align*}
\hat{\beta}_1 &= S_{xy}/S_{xx}, \quad \hat{\beta}_0=\bar{Y}-\hat{\beta}_1\bar{x} \\
R^2 &= 1-\text{SSE}/\text{SST} \\
T &= \hat{\beta}_1/(\hat{\sigma}/\sqrt{S_{xx}}) \sim t(n-2)
\end{align*}
\end{keyformulas}
